\section{Exact dimensions and horizontal components of original cuts}
All rings, ideals, localizations and quotient maps below are the actual
original ones. Dimension is Krull dimension. These statements supply
the local dimension and horizontal-component arguments for a subsequent
scheme cut; no selected component, dimension equality or family is an input.

\begin{lemma}[The whole nonzero principal quotient drops dimension by one]
\label{lem:pl-gauss-affine-principal-quotient-dimension}
\uses{algebra:affine-domain-maximal-height}
\lean{EnvelopeProperty.PLGaussVerticalCurveSurfaceSelectionDimension.principal_quotient_dimension_add_one}
\leanok
Let $A$ be an actual finite-type domain over a field $k$, and let
$0\ne f\in A$ generate a proper ideal. Then
\[
 \dim A/(f)+1=\dim A.
\]
\end{lemma}
\begin{proof}
Since $f$ is a nonzero divisor, the genuine principal-quotient dimension
inequality gives $\dim A/(f)+1\le\dim A$. Choose an actual maximal
ideal $M$ containing the proper ideal $(f)$. Its height equals $\dim A$
by the original finite-type domain maximal-height theorem. The genuine
height inequality for a prime containing $f$ gives
$\operatorname{ht}M\le\dim A/(f)+1$. Substituting this actual maximal
height gives the reverse inequality, proving equality.
\end{proof}

\begin{lemma}[A chosen original minimal prime is isolated on a principal open]
\label{lem:pl-gauss-minimal-prime-isolated-away}
\uses{thm:actual-finite-prime-separator}
\lean{EnvelopeProperty.PLGaussVerticalCurveSurfaceSelectionDimension.exists_away_radical_map_eq}
\leanok
Let $A$ be Noetherian, $I$ an actual ideal, and $P$ an actual minimal
prime over $I$. There is an actual $g\notin P$ such that, for the
unchanged map $\phi:A\to A_g$,
\[
 \sqrt{IA_g}=PA_g.
\]
\end{lemma}
\begin{proof}
The actual minimal primes over $I$ form a finite set. Distinct such
primes cannot contain one another: minimality would force equality.
The original finite prime separator therefore constructs an actual
$g\notin P$ lying in every other minimal prime. The powers of $g$ are
disjoint from $P$, so its actual image $PA_g$ is prime and contracts to
$P$. Localization carries precisely the minimal primes over $I$ that
avoid $g$ to the minimal primes over $IA_g$. Every other original
minimal prime contains $g$ and is excluded; thus the complete set of
minimal primes over $IA_g$ is the singleton $\{PA_g\}$. The radical
of an ideal is the intersection of all its minimal primes. The asserted
equality follows from this actual singleton, without replacing a
supplied component or assuming its localization is irreducible.
\end{proof}

\begin{theorem}[Every original hypersurface component has the exact dimension]
\label{thm:pl-gauss-affine-principal-component-dimension}
\uses{lem:pl-gauss-affine-principal-quotient-dimension,lem:pl-gauss-minimal-prime-isolated-away,algebra:affine-domain-dimension-comparison}
\lean{EnvelopeProperty.PLGaussVerticalCurveSurfaceSelectionDimension.principal_component_dimension_add_one}
\leanok
Let $A$ be an actual finite-type domain over a field $k$, let $f\ne0$,
and let $P$ be ANY actual minimal prime over $(f)$. Then
\[
 \dim A/P+1=\dim A.
\]
There is no incoming height, equidimensionality or dimension certificate.
\end{theorem}
\begin{proof}
Use the preceding lemma to construct $g\notin P$ with
$\sqrt{(f)A_g}=PA_g$. In particular $g\ne0$. The genuine nonzero
principal localization $A_g$ is a finite-type domain over $k$ and has
the entire original dimension $\dim A$. Its image of $f$ remains
nonzero, since the localization map is injective. It is not a unit,
since it belongs to the proper actual prime $PA_g$. The first lemma
therefore gives
$\dim A_g/(f)+1=\dim A_g$. The radical equality says that the two
actual closed prime spectra of $A_g/(f)$ and $A_g/PA_g$ are the same
ordered prime space, so these quotient rings have the same dimension.

The unchanged quotient of the original localization map makes
$A_g/PA_g$ a genuine localization of the actual domain $A/P$ at the
nonzero image of $g$. The canonical localization universal property
identifies it with $(A/P)_{\bar g}$. Nonzero principal localization of
the original finite-type field domain $A/P$ preserves its whole
dimension. Substitute both true dimension identifications in the
preceding equality to obtain $\dim A/P+1=\dim A$. No assertion that
all localizations preserve dimension or that a principal ideal is
prime is used.
\end{proof}

\begin{lemma}[Avoiding central components excludes vertical cut components]
\label{lem:pl-gauss-central-parameter-avoids-cut-components}
\lean{EnvelopeProperty.PLGaussVerticalCurveSurfaceSelectionDimension.parameter_notMem_cut_minimalPrime}
\leanok
Let $A$ be a Noetherian domain, let $t\ne0$, and suppose the actual
element $f$ belongs to NONE of the actual minimal primes over $(t)$.
Then $t$ belongs to NONE of the actual minimal primes over $(f)$.
\end{lemma}
\begin{proof}
Fix an actual $P$ minimal over $(f)$. The principal height theorem gives
$\operatorname{ht}P\le1$. If $t\in P$, then $(t)\subseteq P$.
Since $t$ is a nonzero divisor, the principal ideal $(t)$ has height
at least one. Thus $\operatorname{ht}P\le\operatorname{ht}(t)$.
A prime containing an ideal whose height is at least the prime's
height is minimal over that ideal, so $P$ is an actual minimal prime
over $(t)$. But $f\in P$, contrary to the assumed avoidance. This
proves the assertion directly from the actual principal ideals. It
does not assume that the entire Cartier cut is flat or that the central
fibre is reduced, regular or Cohen--Macaulay.
\end{proof}

For a later induction one must still identify these original affine
component ideals with the actual scheme components containing the
prescribed central curve, retain its original closed inclusion, choose
the actual horizontal component and normalize it, and track the actual
finite dominant lifted central curve and line pullback triangles.
Every selected component is normalized before the next induction step.
The local component producers above do not assert that this global
surface-selection induction has already been assembled.
